%% La coupe A-A : plan de coupe repéré sur la vue de dessus, hachures sur la vue en coupe.
%% La règle des hachures (ce qui est coupé) est la réponse : elle passe par \rappel.
\documentclass[tikz,border=4pt]{standalone}
\usepackage{liaisons}
\begin{document}
\begin{tikzpicture}[x=1cm,y=1cm,
  trait/.style={line width=1.1pt, draw=solideB, line join=round},
  cache/.style={line width=0.7pt, draw=black!55, dash pattern=on 4pt off 3pt},
  plan/.style={line width=0.9pt, draw=solideA, dash pattern=on 9pt off 3pt on 2pt off 3pt},
  axeT/.style={line width=0.5pt, draw=black!55, dash pattern=on 8pt off 3pt on 2pt off 3pt},
  nom/.style={font=\small\bfseries, text=solideB!80!black, inner sep=3pt}]

%% ================= vue de dessus, avec le plan de coupe =================
\draw[trait] (0,0) rectangle (4,2.4);
\draw[trait] (1.2,1.2) circle (0.45);
\draw[trait] (2.9,1.2) circle (0.45);
\draw[axeT] (0.5,1.2) -- (3.6,1.2);
\draw[axeT] (1.2,0.5) -- (1.2,1.9);
\draw[axeT] (2.9,0.5) -- (2.9,1.9);
%% le plan de coupe
\draw[plan] (-0.45,1.2) -- (4.45,1.2);
\draw[-{Stealth[length=7pt,width=6pt]}, line width=1.1pt, draw=solideA] (-0.45,1.2) -- (-0.45,1.95);
\draw[-{Stealth[length=7pt,width=6pt]}, line width=1.1pt, draw=solideA] (4.45,1.2) -- (4.45,1.95);
\node[font=\small\bfseries, text=solideA, anchor=south east] at (-0.35,1.98) {A};
\node[font=\small\bfseries, text=solideA, anchor=south west] at (4.35,1.98) {A};
\node[nom, anchor=north] at (2,-0.2) {Vue de dessus};

%% ================= la vue en coupe A-A =================
\begin{scope}[yshift=-4.2cm]
%% la matière coupée, hachurée ; les trous ne sont pas de la matière
\hachures[draw=solideB, line width=1.1pt]{(0,0) rectangle (0.75,1.3)}
\hachures[draw=solideB, line width=1.1pt]{(1.65,0) rectangle (2.45,1.3)}
\hachures[draw=solideB, line width=1.1pt]{(3.35,0) rectangle (4,1.3)}
\draw[axeT] (1.2,-0.4) -- (1.2,1.7);
\draw[axeT] (2.9,-0.4) -- (2.9,1.7);
\node[font=\small\bfseries, text=solideA, anchor=north] at (2,-0.55) {A\,--\,A};
\end{scope}

%% ================= la règle : c'est la réponse =================
\node[anchor=north west, font=\small, align=left, inner sep=3pt] at (-0.5,-6.1) {%
  \rappel{on hachure \textbf{ce que le plan de coupe traverse}, et rien d'autre :\\
  un trou n'est pas de la matière, il n'est jamais hachuré}};
\end{tikzpicture}
\end{document}
